\section{Resultados e Discussão}

Os resultados apresentados nesta seção referem-se exclusivamente às coletas
realizadas a vácuo.

\subsection{Análise das Respostas dos Sensores}

As respostas médias dos sensores MQ variaram entre as classes e entre as
condições de umidade do solo, conforme apresentado na Tabela
\ref{tab:medias_mq_vacuo}.

\begin{table}[htbp]
    \centering
    \caption{Médias dos sinais dos sensores MQ nas coletas a vácuo.}
    \label{tab:medias_mq_vacuo}
    \scriptsize
    \begin{tabular}{llrrrrrr}
        \hline
        \textbf{Solo} & \textbf{Classe} & \textbf{MQ2} & \textbf{MQ3} &
        \textbf{MQ7} & \textbf{MQ8} & \textbf{MQ135} & \textbf{MQ138} \\
        \hline
        Seco  & Com nematoide  & 8.563,16  & 16.596,73 & 877,35   & 3.818,48 & 19.343,26 & 9.600,27  \\
        Seco  & Sem nematoide  & 15.020,93 & 14.428,36 & 8.215,06 & 6.438,71 & 16.538,29 & 14.843,42 \\
        Úmido & Com nematoide  & 11.395,97 & 14.653,63 & 2.971,30 & 4.222,65 & 17.230,14 & 10.453,63 \\
        Úmido & Sem nematoide  & 10.259,75 & 16.964,65 & 796,36   & 3.390,74 & 20.011,21 & 9.325,97  \\
        \hline
    \end{tabular}

    \vspace{0.1cm}
    {\footnotesize Fonte: Elaborado pelo autor (2026).}
\end{table}

No solo seco, as maiores diferenças entre as classes ocorreram nos sensores
MQ7, MQ2, MQ138 e MQ8. As plantas sem nematoide apresentaram médias superiores
nesses quatro sensores, enquanto MQ3 e MQ135 apresentaram médias superiores nas
plantas com nematoide. Portanto, a separação entre as classes resultou da
resposta conjunta dos sensores, e não do aumento uniforme de todos os sinais.

No solo úmido, parte dessas relações foi invertida e as diferenças entre as
classes foram menos consistentes. Essa alteração do perfil dos sinais indica
maior dificuldade de discriminação nessa condição, o que também foi observado
nos resultados dos classificadores.

\subsection{Resultados da Coleta a Vácuo em Solo Seco}

No solo seco, a X-TCN apresentou o melhor desempenho nas três métricas, como
mostra a Tabela~\ref{tab:resultados_solo_seco_vacuo}.

\begin{table}[htbp]
    \centering
    \caption{Desempenho dos classificadores no solo seco com coleta a vácuo.}
    \label{tab:resultados_solo_seco_vacuo}
    \begin{tabular}{lrrr}
        \hline
        \textbf{Modelo} & \textbf{Acurácia (\%)} &
        \textbf{Acurácia balanceada (\%)} & \textbf{F1 macro (\%)} \\
        \hline
        a) X-TCN         & \textbf{98,47} & \textbf{98,03} & \textbf{98,38} \\
        b) Random Forest & 96,15          & 95,06          & 95,87          \\
        \hline
    \end{tabular}

    \vspace{0.1cm}
    {\footnotesize Fonte: Elaborado pelo autor (2026).}
\end{table}

Em relação ao Random Forest, a X-TCN apresentou ganhos de 2,32 pontos
percentuais em acurácia, 2,97 pontos percentuais em acurácia balanceada e 2,51
pontos percentuais em F1 macro. A rede classificou corretamente todas as
leituras sem nematoide e alcançou recall de 96,07\% nas leituras com nematoide.
O Random Forest também reconheceu todas as leituras sem nematoide, mas obteve
recall de 90,12\% para a classe com nematoide.

A Figura~\ref{fig:matrizes_confusao_solo_seco_vacuo} apresenta as matrizes de
confusão dos dois classificadores no solo seco.

\begin{figure}[htbp]
    \centering
    \includegraphics[width=\linewidth]{figures/matrizes_solo_seco.png}
    \caption{Matrizes de confusão no teste em solo seco: (a) X-TCN e
    (b) Random Forest. As linhas representam as classes reais e as colunas,
    as classes previstas; os percentuais são relativos à classe real.}
    \label{fig:matrizes_confusao_solo_seco_vacuo}

    \vspace{0.1cm}
    {\footnotesize Fonte: Elaborado pelos autores (2026).}
\end{figure}

A X-TCN reduziu o número de erros de 643 para 256 leituras, correspondendo a
uma redução de 60,19\%. Esse resultado mostra que a representação temporal dos
sinais favoreceu a classificação no solo seco. Contudo, como as leituras de
uma mesma coleta são correlacionadas, o desempenho deve ser confirmado em
novas coletas independentes.

\subsection{Resultados da Coleta a Vácuo em Solo Úmido}

No solo úmido, o Random Forest superou a X-TCN nas três métricas avaliadas,
conforme a Tabela~\ref{tab:resultados_solo_umido_vacuo}.

\begin{table}[htbp]
    \centering
    \caption{Desempenho dos classificadores no solo úmido com coleta a vácuo.}
    \label{tab:resultados_solo_umido_vacuo}
    \begin{tabular}{lrrr}
        \hline
        \textbf{Modelo} & \textbf{Acurácia (\%)} &
        \textbf{Acurácia balanceada (\%)} & \textbf{F1 macro (\%)} \\
        \hline
        Random Forest & \textbf{74,46} & \textbf{62,96} & \textbf{63,33} \\
        X-TCN         & 63,58          & 41,52          & 38,87          \\
        \hline
    \end{tabular}

    \vspace{0.1cm}
    {\footnotesize Fonte: Elaborado pelo autor (2026).}
\end{table}

O Random Forest obteve recall de 84,61\% para as leituras com nematoide e de
41,31\% para as leituras sem nematoide. A diferença entre a acurácia de 74,46\%
e a acurácia balanceada de 62,96\% mostra que o resultado foi influenciado pelo
desbalanceamento entre as classes.

A X-TCN alcançou recall de 83,04\% para a classe com nematoide, mas não
reconheceu nenhuma leitura da classe sem nematoide. Por esse motivo, apesar da
acurácia de 63,58\%, a acurácia balanceada caiu para 41,52\% e o F1 macro para
38,87\%. O resultado mostra que o padrão temporal identificado no solo seco
não se manteve estável no solo úmido.

A Figura~\ref{fig:matrizes_confusao_solo_umido_vacuo} apresenta as matrizes de
confusão dos dois classificadores no solo úmido.

\begin{figure}[htbp]
    \centering
    \includegraphics[width=\linewidth]{figures/matrizes_solo_umido.png}
    \caption{Matrizes de confusão no teste em solo úmido: (a) X-TCN e
    (b) Random Forest. As linhas representam as classes reais e as colunas,
    as classes previstas; os percentuais são relativos à classe real.}
    \label{fig:matrizes_confusao_solo_umido_vacuo}

    \vspace{0.1cm}
    {\footnotesize Fonte: Elaborado pelos autores (2026).}
\end{figure}

\subsection{Comparação entre as Condições de Umidade do Solo}

Os dois modelos apresentaram desempenho superior no solo seco. As diferenças
entre as condições são apresentadas na Tabela
\ref{tab:comparacao_umidade_vacuo}.

\begin{table}[htbp]
    \centering
    \caption{Comparação do desempenho entre solo seco e solo úmido nas coletas a vácuo.}
    \label{tab:comparacao_umidade_vacuo}
    \footnotesize
    \begin{tabular}{llrrr}
        \hline
        \textbf{Modelo} & \textbf{Condição} & \textbf{Acurácia (\%)} &
        \textbf{Acurácia bal. (\%)} & \textbf{F1 macro (\%)} \\
        \hline
        Random Forest & Solo seco  & 96,15 & 95,06 & 95,87 \\
                      & Solo úmido & 74,46 & 62,96 & 63,33 \\
                      & Diferença  & 21,69 & 32,10 & 32,54 \\
        \hline
        X-TCN         & Solo seco  & 98,47 & 98,03 & 98,38 \\
                      & Solo úmido & 63,58 & 41,52 & 38,87 \\
                      & Diferença  & 34,89 & 56,51 & 59,51 \\
        \hline
    \end{tabular}

    \vspace{0.1cm}
    {\footnotesize Fonte: Elaborado pelo autor (2026).}
\end{table}

No Random Forest, a mudança do solo seco para o solo úmido reduziu a acurácia
em 21,69 pontos percentuais, a acurácia balanceada em 32,10 pontos percentuais
e o F1 macro em 32,54 pontos percentuais. Na X-TCN, as reduções foram de 34,89,
56,51 e 59,51 pontos percentuais, respectivamente.

A queda mais acentuada das métricas balanceadas confirma que a condição de
umidade afetou principalmente a distinção equilibrada entre as classes. A
X-TCN foi mais sensível a essa mudança, enquanto o Random Forest manteve maior
capacidade de classificação no solo úmido.

\subsection{Comparação dos Modelos de Classificação}

A superioridade dos modelos dependeu da condição do solo. No solo seco, a
X-TCN apresentou os melhores resultados, alcançando 98,47\% de acurácia,
98,03\% de acurácia balanceada e 98,38\% de F1 macro. No solo úmido, o Random
Forest foi superior, com 74,46\%, 62,96\% e 63,33\%, respectivamente.

Ao considerar conjuntamente todas as leituras coletadas a vácuo, o Random
Forest atingiu 90,83\% de acurácia e a X-TCN atingiu 89,91\%. Essa diferença
global é pequena, mas oculta comportamentos distintos: a X-TCN foi mais eficaz
no solo seco e falhou na identificação da classe sem nematoide no solo úmido,
enquanto o Random Forest apresentou maior estabilidade entre as condições.

Assim, os resultados não indicam superioridade geral de um único modelo. A
X-TCN mostrou maior capacidade de explorar os padrões temporais quando as
respostas dos sensores estavam mais bem separadas. O Random Forest, por sua
vez, foi menos afetado pela alteração do perfil dos sinais no solo úmido.

\subsection{Importância das Variáveis e dos Sensores}

% Nomenclatura provisória solicitada: o modelo arbóreo executado é ExtraTreesClassifier.
Para comparar a contribuição das entradas nos dois classificadores, foi
realizada uma análise adicional com os mesmos conjuntos de treinamento e teste,
separados por coleta. Em cada classificador foram ajustadas duas versões: uma
com os seis sinais MQ brutos (MQ2, MQ3, MQ7, MQ8, MQ135 e MQ138) e outra com
esses sinais acrescidos do índice de umidade do solo, da temperatura e da
pressão interna. Os sinais MQ brutos foram mantidos nas duas versões porque os
MQ corrigidos utilizados na análise anterior já incorporam uma compensação
estatística baseada nas variáveis ambientais. Assim, os resultados desta
ablação não substituem os percentuais apresentados nas tabelas anteriores.
O modelo arbóreo utilizou 900 árvores, seleção de atributos pela raiz
quadrada do número de entradas e mínimo de dez amostras por folha. A X-TCN
foi treinada por 60 épocas com janelas de 128 leituras e passo de 64; a
semente aleatória foi 42 nos dois casos.

O teste incluiu 22.115 leituras de 12 coletas, das quais 16.692 leituras de
oito coletas foram classificadas operacionalmente como solo seco e 5.423
leituras de quatro coletas como solo úmido. Cada modelo foi treinado uma única
vez com todas as coletas de treinamento; a separação por condição do solo foi
aplicada apenas à avaliação. Na X-TCN, a probabilidade atribuída a cada
leitura corresponde à média das probabilidades das janelas temporais que a
contêm. A Tabela~\ref{tab:ablacao_ambiente_vacuo} apresenta os resultados
obtidos por leitura.

\begin{table}[htbp]
    \centering
    \caption{Desempenho dos classificadores com e sem variáveis ambientais nas coletas a vácuo.}
    \label{tab:ablacao_ambiente_vacuo}
    \scriptsize
    \begin{tabular}{llrrrr}
        \hline
        \textbf{Modelo} & \textbf{Entradas} &
        \textbf{Acurácia total (\%)} &
        \textbf{Bal. total (\%)} &
        \textbf{Bal. seco (\%)} &
        \textbf{Bal. úmido (\%)} \\
        \hline
        Random Forest & MQ              & 89,86 & 89,90 & 99,93 & 38,56 \\
        Random Forest & MQ + ambiente   & 87,93 & 87,89 & 99,92 & 33,30 \\
        X-TCN         & MQ              & 80,06 & 79,77 & 91,50 & 26,07 \\
        X-TCN         & MQ + ambiente   & 89,41 & 89,54 & 95,12 & 46,15 \\
        \hline
    \end{tabular}

    \vspace{0.1cm}
    {\footnotesize Bal. = acurácia balanceada. Fonte: Elaborado pelo autor (2026).}
\end{table}

Para medir a importância das entradas nos dois modelos, foi adotada a
\emph{importância por permutação} (\emph{permutation feature importance}),
conforme o princípio de perturbar uma variável por vez e medir a alteração do
desempenho do modelo \cite{fisher2019modelreliance}. Em cada repetição, os
valores de uma coluna foram embaralhados entre as leituras do teste completo,
mantendo-se o modelo treinado e as demais colunas. A importância corresponde à
diferença entre a acurácia balanceada original e a obtida após a permutação,
expressa em pontos percentuais. Foram realizadas três permutações por variável
em cada cenário. O mesmo procedimento e a mesma métrica foram usados para o
Random Forest e para a X-TCN; nesta última, as janelas foram reconstruídas a
partir das leituras perturbadas antes da predição. A
Tabela~\ref{tab:importancia_variaveis_vacuo} apresenta a média das três quedas.

\begin{table}[htbp]
    \centering
    \caption{Importância por permutação no teste completo: queda da acurácia balanceada, em pontos percentuais.}
    \label{tab:importancia_variaveis_vacuo}
    \scriptsize
    \begin{tabular}{lrrrr}
        \hline
        \textbf{Variável} &
        \textbf{RF: MQ} &
        \textbf{X-TCN: MQ} &
        \textbf{RF: MQ + amb.} &
        \textbf{X-TCN: MQ + amb.} \\
        \hline
        Pressão interna & --    & --    & 7,13  & 23,49 \\
        Umidade do solo & --    & --    & -0,37 & 8,12  \\
        Temperatura     & --    & --    & -2,60 & -1,01 \\
        MQ2             & 2,40  & -3,68 & -0,24 & 17,64 \\
        MQ3             & 3,84  & -8,92 & 1,13  & -3,77 \\
        MQ7             & 3,88  & -1,14 & 1,63  & 3,07  \\
        MQ8             & 9,28  & -5,11 & 1,96  & 5,25  \\
        MQ135           & 9,40  & 8,20  & 0,57  & 3,77  \\
        MQ138           & 0,17  & -5,13 & 1,43  & 10,59 \\
        \hline
    \end{tabular}

    \vspace{0.1cm}
    {\footnotesize -- = variável ausente do modelo. Fonte: Elaborado pelo autor (2026).}
\end{table}

Com as variáveis ambientais, a pressão interna produziu a maior queda de
acurácia balanceada nos dois modelos: 7,13 pontos percentuais no Random Forest
e 23,49 na X-TCN. Na X-TCN, também se destacaram MQ2 (17,64), MQ138 (10,59)
e o índice de umidade do solo (8,12 pontos percentuais). Sem as entradas
ambientais, MQ135 apresentou a maior queda nos dois classificadores. Essas
diferenças mostram que a posição de um sensor no ranking depende das demais
entradas oferecidas ao modelo. Valores negativos indicam que, nessa avaliação,
o embaralhamento da variável aumentou a acurácia balanceada; não representam
uma influência física negativa do sensor sobre a presença de nematoides.

A inclusão das variáveis ambientais elevou a acurácia balanceada total da
X-TCN de 79,77\% para 89,54\%, mas reduziu a do Random Forest de 89,90\% para
87,89\%. No recorte úmido, a X-TCN passou de 26,07\% para 46,15\% e o Random
Forest passou de 38,56\% para 33,30\%. Portanto, o ganho da X-TCN com as
informações ambientais não eliminou a dificuldade de classificação nessa
condição. A distinção entre solo seco e úmido foi feita por um índice
operacional do sensor de solo, sem medição calibrada da água em diferentes
profundidades. Como o teste úmido contém apenas quatro coletas e a condição de
umidade também varia com o dia de coleta, esses resultados não isolam um efeito
causal da água sobre os gases ou sobre os nematoides. A importância preditiva
da pressão também pode refletir características da câmara e da operação do
sistema de vácuo.
